% subsection: 階差等比複合型 \, $\displaystyle a_{n+1}=pa_n+f(n)$

\subsection{\texorpdfstring{階差等比複合型 \, $\displaystyle a_{n+1}=pa_n+f(n)$}{階差等比複合型}}\label{節：階差等比複合型}
  \begin{bookpoint}
    \textbf{同じ更新をする量を引く}\par
    付加項を消すことを目標に,引く式の形と係数を決める.候補を見つける計算と,候補が元の式を満たす確認を区別しよう.
  \end{bookpoint}
  \subsubsection{\texorpdfstring{多項式型 \, $\displaystyle a_{n+1}=pa_n+f(n)$}{多項式型}}
    \bookstep{式の形から、追う量を選ぶ}
    \sectionref*{節：特性方程式型}では,同じ更新をする定数を引いた.
    今回は付加項$f(n)$が$n$によって変わるため,引く量も$n$によって変えてみる.
    \sectionref*{節：補助数列の設計}で述べた方針に従い,同じ更新をする既知の量を探そう.
    本文では$f(n)$が多項式の場合を考える.$p=1$なら\bookterm{difference}{階差}型として和を求められるので,まず$p\ne1$とする.

    \bookstep{引く量を選ぶ}
    $b_n=a_n-g(n)$を等比型にするための条件は
    \[
      g(n+1)=pg(n)+f(n)
    \]
    である.これは,\,$g(n)$自身にも同じ更新をさせるという条件である.
    条件が満たされれば,
    \[
      \begin{array}{rrl}
        &a_{n+1}={}&pa_n+f(n)\\
        -\,)&g(n+1)={}&pg(n)+f(n)\\ \hline
        &a_{n+1}-g(n+1)={}&p\{a_n-g(n)\}
      \end{array}
    \]
    となり,付加項が消える.
    \sectionref*{節：非同次と解の組合せ}で見た「一つの既知の解と\bookterm{homogeneous}{同次}の解に分ける」考え方が,
    ここでは$a_n=g(n)+b_n$として現れている.
    $g(n)$を見つける計算と,残る\bookterm{homogeneous}{同次}の解$b_n$を\bookterm{initial}{初期条件}から選ぶ計算を分けているのである.

    $f(n)$が$m$次式なら,\,$p\ne1$のときは$m$次の多項式
    \[
      g(n)=x_0+x_1n+\cdots+x_mn^m
    \]
    を候補にする.$g(n+1)-pg(n)$の最高次の係数は$(1-p)x_m$である.
    これを$f(n)$の最高次の係数に合わせ,次に一段低い係数を合わせる.
    $1-p\ne0$なので,この順に$x_m,x_{m-1},\ldots,x_0$を決められる.
    未知の係数を比較して求めるこの方法を\bookterm{undetermined}{未定係数法}という.

    例えば$a_{n+1}=2a_n+n-1$では,\,$g(n)=un+v$と置くと
    \[
      u=2u+1,\qquad u+v=2v-1
    \]
    となるので,\,$g(n)=-n$である.
    従って$a_n+n$が2倍される量となる.
    この例の計算は\bookproblemref{practice-basic}{4}{標準問4}で確認する.

    \bookstep{変化を求めて、元の項へ戻す}
    $b_{n+1}=pb_n$を解き,\,$a_n=b_n+g(n)$へ戻す.
    $p\ne0,1$なら
    \[
      a_n=\{a_1-g(1)\}p^{n-1}+g(n)
    \]
    となる.$p=0$なら直接$n\ge2$で$a_n=f(n-1)$と求める.
    一方,\,$p=1$で$f(n)$が$m$次なら,付加項を足し戻す多項式は通常$m+1$次になる.
    例えば$a_{n+1}=a_n+n$では,\,$g(n)=\frac{n(n-1)}2$が同じ更新を満たす.

  \subsubsection{\texorpdfstring{指数型  \, $\displaystyle a_{n+1}=pa_n+q^n$}{指数型}}
    前節では$f(n)$が多項式である場合を扱った.
    ここでは$f(n)$が特に指数関数である場合を扱う.
    指数関数である場合にはもう一つ別の解法が存在する.
    \begin{align*}
      a_{n+1}=pa_n+q^n
    \end{align*}
    以下では$q\ne0$とする($q=0$なら$n\ge1$で$q^n=0$なので等比型になる).
    両辺を$q^n$で割って,
    \begin{align*}
      &\frac{a_{n+1}}{q^n}=\frac{p}{q}\cdot\frac{a_n}{q^{n-1}}+1\\
    \end{align*}
    ここで,\,$b_n=\frac{a_n}{q^{n-1}}$と置くと
    \begin{align*}
      &b_{n+1}=\frac{p}{q}b_n+1
    \end{align*}
    あとは\bookterm{characteristic}{特性方程式}型の\bookterm{recurrence}{漸化式}として解くことになる.
    ただし$p=q$のときは$b_{n+1}=b_n+1$(\bookterm{arithmetic}{等差数列})になり,\bookterm{fixed}{不動点}$\alpha$が存在しない(\sectionref{節：特性方程式型}の例題(3)と同じ状況である).
    この場合は$b_n=b_1+(n-1)$として直接求まる.
    例えば$a_1=1,\,a_{n+1}=2a_n+2^n$(\,$p=q=2$)ならば$b_n=\dfrac{a_n}{2^{n-1}}$が\bookterm{arithmetic}{等差数列}となり,\,$a_n=n\cdot2^{n-1}$が得られる.
